\section*{Chapter \ref{ch:g11:asymmetric-information} --- Asymmetric Information}
% figdata: figdata/grade-11/08-asymmetric-information/out/ (hie, lemons, lemons_end); Brenland numbers invented.

\begin{solution}{exo:g11:asymmetric-information:1}
\omterm{def:g11:asymmetric-information:adverse-selection}{Adverse selection} (a hidden characteristic before the deal); \omterm{def:g11:asymmetric-information:moral-hazard}{moral hazard} (a hidden action after it); \omterm{def:g11:asymmetric-information:adverse-selection}{adverse selection}.
\end{solution}

\begin{solution}{exo:g11:asymmetric-information:2}
\omterm{def:g11:asymmetric-information:moral-hazard}{Moral hazard} (a principal--agent case); \omterm{def:g11:asymmetric-information:adverse-selection}{adverse selection}; \omterm{def:g11:asymmetric-information:moral-hazard}{moral hazard}.
\end{solution}

\begin{solution}{exo:g11:asymmetric-information:3}
\omterm{def:g11:asymmetric-information:moral-hazard}{Moral hazard}; \omterm{def:g11:asymmetric-information:adverse-selection}{adverse selection} (the high rate attracts the riskier borrowers); \omterm{def:g11:asymmetric-information:moral-hazard}{moral hazard}.
\end{solution}

\begin{solution}{exo:g11:asymmetric-information:4}
$(2170 - 1325)/2170 \approx 39\,\%$ less.
\end{solution}

\begin{solution}{exo:g11:asymmetric-information:5}
From 93.1\,\% to 76.1\,\%: 17 points. The two plans' estimates are uncertain (the 50\,\% plan had few people, and the margin of error on its difference from free care is
large), so a gap of about 110 dollars between them may be chance.
\end{solution}

\begin{solution}{exo:g11:asymmetric-information:6}
A car of unknown quality is worth $0.5 \times 10\,000 + 0.5 \times 3000 = 6500$ to a buyer, less than the 8000 a good car's owner wants: good cars are withdrawn. Buyers,
knowing that, value a car offered at 3000; only bad cars are sold, at a price between 2000 and 3000.
\end{solution}

\begin{solution}{exo:g11:asymmetric-information:7}
A car of unknown quality is now worth $0.8 \times 10\,000 + 0.2 \times 3000 = 8600$ to a buyer, more than 8000: all cars are sold, at a price between 8000 and 8600. Bad cars'
owners gain most.
\end{solution}

\begin{solution}{exo:g11:asymmetric-information:8}
A good car's seller gets $10\,000 - 200 = 9800$ with a warranty, more than the 8000 he values his car at: he offers one. A bad car's seller would get $10\,000 - 8000 = 2000$,
less than the 3000 he can get by selling without a warranty as a known bad car: he does not offer one. Buyers can trust the warranty as a sign of a good car.
\end{solution}

\begin{solution}{exo:g11:asymmetric-information:9}
The insurer pays $2000 - 500 = 1500$ when the loss strikes: $0.1 \times 1500 = 150$ a year expected. The insured person pays 500 when it strikes: 50 a year expected.
\end{solution}

\begin{solution}{exo:g11:asymmetric-information:10}
The able worker must gain by taking it, $30 - c_H \ge 20$, and the other must not, $30 - c_L < 20$: $c_H \le 10 < c_L$.
\end{solution}

\begin{solution}{exo:g11:asymmetric-information:11}
Qualities from 2000: about 8800 in round 0 and 5900 in round 15. Qualities from 0: about 7300 in round 0 and about 200 in round 15. About half the cars are traded at the
end, as in the example, where the price is 6000 and half the cars are sold; the small differences come from the random draw of 1000 cars.
\end{solution}

\begin{solution}{exo:g11:asymmetric-information:12}
The model predicts that the trucks offered for sale are the worse ones, so trucks bought used should need more maintenance than trucks of the same age and mileage that never
changed hands. The study (Bond 1982) found no such difference. Possible reasons: buyers of trucks, often tradesmen, can inspect a truck or have it inspected; dealers with a reputation
stand behind what they sell; and a truck's owner may sell it for reasons unrelated to its quality, such as a change of business.
\end{solution}

\begin{solution}{exo:g11:asymmetric-information:13}
All three insured: premium $(100 + 300 + 800)/3 = 400$; the low risks will pay at most 150 and leave. The other two: premium $(300 + 800)/2 = 550$; the middle risks will pay
at most 350 and leave. The high risks alone: premium 800, which they accept (up to 850). Only the highest risks remain insured.
\end{solution}

\begin{solution}{exo:g11:asymmetric-information:14}
On a fixed salary his pay does not depend on profit, so he does not make the effort. With a share $s$ he gains $s \times 600$ from the effort, which must cover its cost:
$600s \ge 200$, $s \ge 1/3$. The owners give up a third of the profit, but keep $\tfrac23 \times 1600 \approx 1067$ before the salary, more than 1000. The new problems: the
manager now bears the risk of profits that depend on luck, and will ask for more pay in return; and he may raise measured profit in ways that harm the firm later.
\end{solution}

\begin{solution}{exo:g11:asymmetric-information:15}
The premium is $(100 + 300 + 800)/3 = 400$. The high risks gain: they paid 800 in the voluntary market and now pay 400. The middle risks value full cover at
350 and must pay 400: they lose 50 against staying uninsured. The low risks value it at 150 and pay 400: they lose 250. Compulsory cover keeps everyone insured by moving
part of the high risks' cost onto the others.
\end{solution}

\begin{solution}{pb:g11:asymmetric-information:1}
\textbf{1.} All of them, since $1.5q > q$ for every car.

\textbf{2.} $0.5q$ on average, $0.5 \times 6000 = 3000$ dalens.

\textbf{3.} The cars with $q \le p$.

\textbf{4.} $(2000 + p)/2$.

\textbf{5.} $1.5 \times (2000 + p)/2 = 1500 + 0.75p$.

\textbf{6.} Buyers pay $p$ only if $1500 + 0.75p \ge p$, that is $p \le 6000$; with many buyers, the price is 6000.

\textbf{7.} $(6000 - 2000)/8000 = 50\,\%$.

\textbf{8.} The average quality traded is 4000, worth 6000 to a buyer, which is the price: the buyers' average gain is zero.

\textbf{9.} A car traded brings $1.5 \times 4000 - 4000 = 2000$ on average; half the cars are traded, so 1000 per car in the town, a third of the 3000 of question 2.

\textbf{10.} The average car offered is worth $1.5 \times p/2 = 0.75p < p$ at every positive price: the price falls to zero and the market vanishes.

\textbf{11.} A seller gains $1.5q - 400$ by inspecting and selling, against $q$ by keeping his car: he inspects when $0.5q \ge 400$, $q \ge 800$, so every seller inspects.

\textbf{12.} All the cars; the average price is $1.5 \times 6000 = 9000$.

\textbf{13.} $0.5 \times 6000 - 400 = 2600$ dalens.

\textbf{14.} By $2600 - 1000 = 1600$ dalens per car.

\textbf{15.} A warranty costs the seller of a bad car more than the seller of a good one, so only good cars carry one: like the inspection, it tells buyers the quality.

\textbf{16.} $1.2 \times (2000 + p)/2 = p$: $1200 + 0.6p = p$, $p = 3000$; the share traded is $(3000 - 2000)/8000 = 12.5\,\%$.

\textbf{17.} The average car offered is worth $2 \times (2000 + p)/2 = 2000 + p > p$ at every price, so buyers keep bidding until every car is offered: all cars are traded.

\textbf{18.} The gap between buyers' and sellers' values is what holds the market together; when it is wide, buyers will pay for an average car enough to keep the good
ones in, and when it is narrow, the good cars leave.

\textbf{19.} He trades again and again with the town; a buyer cheated tells others, and the dealer loses future sales worth more than the gain from one bad car, as in
the \omterm{def:g11:oligopoly-and-games:repeated-game}{repeated game} of \cref{ch:g11:oligopoly-and-games}.

\textbf{20.} Without inspection the price settles at 6000 dalens and half the cars are traded; with an inspection costing 400, every car is traded, at an average price of 9000.
\end{solution}
